\section{Experimental Setup}
\label{sec:setup}

We design experiments to answer four research questions, each corresponding to a contribution from the
Introduction:

\textbf{RQ1 (Existence):} Is generalization gap learnable from weight tensors at Spearman $r > 0.5$, and is
gap genuinely independent of test accuracy (A1 audit)?

\textbf{RQ2 (Architecture ranking):} Which encoder architecture achieves the highest gap Spearman, and with
what statistical confidence?

\textbf{RQ3 (Gap-specific signal):} Do gap predictions from the best encoder contain information about true
gap that is statistically independent of test accuracy?

\textbf{RQ4 (Differential advantage):} Do equivariant encoders show larger Spearman improvement on gap than
on test accuracy ($\Delta > 0.02$) --- the target-specificity mechanism hypothesis?

\subsection{Dataset}

\textbf{Unterthiner CIFAR-10 CNN Zoo.} We evaluate on the model zoo from \citet{unterthiner2020}: $N = 10{,}000$
small convolutional neural networks trained on CIFAR-10 with a fixed 4-layer architecture and varying
hyperparameters (learning rate $\in \{0.0001, 0.001, 0.01, 0.1\}$, weight decay, optimizer $\in \{\text{SGD, Adam,
RMSprop}\}$). Each model provides a weight tensor of dimension $D = 33{,}890$ (all parameters concatenated)
and training logs from which we compute:
\begin{itemize}
\item $y^{\text{acc}}$ = test accuracy at the epoch with best validation accuracy
\item $y^{\text{gap}} = y^{\text{train}} - y^{\text{acc}}$ (generalization gap at convergence)
\end{itemize}

\textbf{Rationale.} The Unterthiner zoo is the standard benchmark for weight-space model property prediction,
enabling direct comparison with prior encoders. Its dense coverage of hyperparameter configurations
provides meaningful variance in both test accuracy and generalization gap.

\textbf{Split.} 80\% training (8,000 models), 10\% validation (1,000 models), 10\% test (1,000 models).
Stratified by hyperparameter configuration, seed 42. Test split is never used during encoder training
or hyperparameter selection.

\textbf{A1 Audit.} Before any training, we verify that generalization gap and test accuracy are not collinear:
Spearman(gap, $-$test\_acc) $= -0.142$ ($|-0.142| \ll 0.95$). Gap and test accuracy are only weakly negatively
correlated in this zoo, confirming that gap prediction is a genuinely distinct task.

\subsection{Encoders (Baselines and Proposed)}

We compare four weight-space encoders spanning the range from non-equivariant to fully equivariant:

\begin{table}[h]
\centering
\caption{Encoder architectures evaluated.}
\label{tab:encoders}
\begin{tabular}{@{}llll@{}}
\toprule
Encoder & Equivariance Type & Architecture & Source \\
\midrule
FlatMLP & None (position-indexed) & 3-layer MLP on sorted weights & \citet{unterthiner2020} \\
DWSNet & Within-layer (row/col) & Equivariant linear layers per matrix & \citet{navon2023} \\
NFT & Cross-layer (sequence attention) & Multi-head attention on weight matrix sequence & \citet{zhou2023} \\
GNN & Graph-structured & Message passing over neuron graph & \citet{kofinas2024} \\
\bottomrule
\end{tabular}
\end{table}

\textbf{Why these baselines.} FlatMLP is the foundational baseline from the zoo benchmark paper --- any claim
about gap learnability must be benchmarked against it. DWSNet, NFT, and GNN are the three dominant
equivariant architectures in the literature, each representing a distinct inductive bias: within-layer
averaging, cross-layer attention, and graph connectivity, respectively. Testing all three allows us to
attribute architecture-specific effects rather than ``equivariance in general.''

\textbf{Dual-target design.} For RQ4, all four encoders are trained on each target (gap and test\_acc)
independently, with identical hyperparameter budget and training procedure. This controls for
zoo-specific effects and isolates target-specificity.

\subsection{Training Protocol}

All encoders use:
\begin{itemize}
\item Optimizer: Adam
\item Learning rate search: 3-trial random search over $\{5\times10^{-4}, 1\times10^{-3}, 2\times10^{-3}\}$
\item Batch size: 64
\item Epochs: 100
\item Learning rate schedule: cosine (NFT, GNN); none (FlatMLP, DWSNet)
\item Selection criterion: best validation Spearman across 3 trials
\item Random seed: 42 for all data splits, model initialization, and training
\end{itemize}

\textbf{Computational resources.} All experiments run on a single GPU. NFT and GNN require $\sim$3$\times$ more compute
per trial than FlatMLP due to attention and message-passing overhead.

\textbf{Acknowledged limitation.} The pre-specified protocol called for 50-trial random search. Resource
constraints reduced this to 3 trials. The impact is most evident in FlatMLP test\_acc Spearman ($r = 0.279$
vs.\ literature $\approx 0.85$) and is discussed in Section~\ref{sec:discussion}.

\subsection{Evaluation Metrics}

\textbf{Spearman rank correlation} (primary): $\rho(\hat{y}, y)$ on the test split ($N = 1{,}000$). Chosen
because practitioners care about relative ranking (which models are least overfit?) rather than absolute
predictions.

\textbf{Bootstrap 95\% confidence intervals} (uncertainty quantification): $N_{\text{boot}} = 1{,}000$
resamples, seed 42, percentile CI. Applied to all Spearman estimates and to $\Delta$ values.

\textbf{Differential advantage} (RQ4): $\Delta(e) = [\rho_{\text{gap}}(e) - \rho_{\text{gap}}(\text{FlatMLP})] - [\rho_{\text{acc}}(e) - \rho_{\text{acc}}(\text{FlatMLP})]$ for each equivariant encoder $e$. Gate criterion: $\Delta > 0.02$ for $\geq 2$ of 3 equivariant encoders.

\textbf{Partial Spearman} (RQ3): $\rho_{\text{partial}}(\hat{y}^{\text{gap}}, y^{\text{gap}} | y^{\text{acc}})$ via rank residual regression (described in Section~\ref{sec:method}). Threshold: $\rho_{\text{partial}} > 0$, $p < 0.05$.

\textbf{Mean Squared Error} (secondary): MSE on the test split, reported alongside Spearman.

\subsection{Statistical Validity}

Experiments follow a pre-registered design (Phase 2B verification plan):
\begin{itemize}
\item A1 audit computed before any training
\item Threshold values ($\Delta > 0.02$, $r > 0.5$, partial $\rho > 0$) pre-specified
\item Test split never used for hyperparameter selection
\item Bootstrap CI used for all reported statistical comparisons
\end{itemize}

All gate decisions (PASS/FAIL) are made against pre-specified thresholds, not post-hoc.
